%! Function Composition and Decomposition — Unit 1, Lesson 1.3 — Lesson Plan
% Section order per LESSON_SHAPE.md §5. REPLAN NOTE: transformations are now Lesson 1.5, AFTER
% this lesson, so nothing here treats f(x+h) as prior knowledge — 1.5 appears only as a forward
% promise. Inverse functions are now Lesson 1.4 (the old plan said 1.6); every cross-reference
% has been renumbered. 1.2 (Function Operations) is the lesson immediately before this one.
\documentclass[10pt]{article}
\usepackage{precalculus-article}
\usepackage{precalculus-boxes}
\usepackage{graphicx}
\graphicspath{{images/}}
\newcommand{\TallMath}[1]{$\displaystyle #1\rule[-1.4em]{0pt}{3.2em}$}
\newcommand{\CourseName}{Precalculus}
\newcommand{\MeetingLength}{60 minutes}
\newcommand{\UnitNumberName}{Unit 1: Functions and Change \quad}
\newcommand{\LessonNumberName}{Lesson 1.3: Function Composition and Decomposition}

\begin{document}
\begin{center}
    {\Large\bfseries \CourseName \\
    {\small\bfseries \UnitNumberName \LessonNumberName}}
\end{center}

\begin{tcolorbox}[colback=lilac, colframe=plum, boxrule=0.9pt, arc=2mm,
    left=2mm, right=2mm, top=1.5mm, bottom=1.5mm]
    \textbf{Primary Objective:} Students will read every composition $f(g(x))$
    through the one number that sits between the two functions --- the
    \textit{handoff}, which is the inner function's answer and the outer
    function's question at once; evaluate composites from a table and from
    formulas by naming that number first; build $(f \circ g)(x)$ by substituting
    $g(x)$ for every $x$ in $f(x)$; explain non-commutativity and the domain of a
    composite as facts about the handoff; and decompose a function into an outer
    rule and an inner rule.
    \hfill\textit{\small (CED Topic 2.7 --- Skills 1.C, 2.B)}
\end{tcolorbox}

\boxguard[22]
\begin{skillbox}[Lesson Flow --- Gradual Release (60 minutes)]{lilac}
\small
\renewcommand{\arraystretch}{1.25}
\begin{tabularx}{\linewidth}{@{} l c X @{}}
\textbf{Phase} & \textbf{Min} & \textbf{What students are doing, and where it lives} \\
\midrule
Warm-Up                       & 5  & Silent, individual spiral review --- warm-up sheet. \\
Hook                          & 2  & Whole-class vote and discussion --- hook box atop the guided notes. \\
\textbf{I Do} --- model       & 10 & Watch and annotate while the teacher thinks aloud --- the THE HANDOFF and FINDING IT rows: definitions, displays, problems 1 and 3. \\
\textbf{We Do} --- guided     & 21 & Fill the notes together, every problem cold-called --- problems 2, 4--6, then the BUILDING IT and TWO HANDOFFS rows (7--14; the teacher works 7 and 11). \\
\textbf{You Do} --- independent & 14 & Work alone while the teacher circulates --- the You Do row, problems 15--18. \\
Debrief                       & 5  & Whole-class share-out and the headline. \\
Homework Launch               & 3  & Start homework problems 1 and 2 alone; the teacher reads problem 2 over shoulders. \\
\end{tabularx}

\smallskip
\textbf{The whole release lives in the guided notes} --- there is no activity sheet to raid when
the clock slips. Four ideas share the block (the handoff, tables, building, and both orders plus
decomposition), so the guided phase runs 21 minutes. If you fall behind, cut \textbf{problem 10}
(the second pair of formulas in the BUILDING IT row) first, then problem 14, and protect the
\textbf{TWO HANDOFFS} row: problem 12 is the lesson, and homework problem 2 cannot be done
without it. Never cut the Homework Launch to zero --- it is the only formative read of the period
now that there is no exit ticket.
\end{skillbox}

\renewcommand{\arraystretch}{1.6}
\begin{skillbox}[Priority Ideas \& Skills]{goldbox}
\begin{tabularx}{\linewidth}{@{}X X@{}}
\textbf{Priority Skills} &
\textbf{Key Understandings} \\
\midrule
\textbf{AP Skill 1.C} --- Construct new functions, using compositions, that may be
useful in modeling contexts, criteria, or data. &
\begin{itemize}[leftmargin=1.2em,itemsep=2pt,topsep=0pt]
  \item $(f \circ g)(x) = f(g(x))$: the outputs of $g$ become the inputs of $f$.
        The number that changes hands --- the \textbf{handoff} --- is the whole
        mechanism. (2.7.A.1)
  \item Composition links two quantities no single formula connects; in context
        the handoff is a real quantity with its own units. (2.7.B.1)
  \item Because the handoff must be a value the outer function accepts, the
        domain of a composite is restricted to those $x$ whose $g(x)$ lands in
        $f$'s domain. (2.7.A.1)
\end{itemize} \\
\textbf{AP Skill 2.B} --- Construct equivalent representations of functions and
expressions. &
\begin{itemize}[leftmargin=1.2em,itemsep=2pt,topsep=0pt]
  \item A composite value can be found from a table, a graph, or an analytic
        expression --- always by naming the handoff first, and all three routes
        agree. (2.7.A.2)
  \item Build $f(g(x))$ analytically by substituting $g(x)$ for \textit{every}
        $x$ in $f(x)$; the handoff is now an expression, not a number. (2.7.B.2)
  \item \textbf{Decomposition} is composition read backwards: name what you
        would compute first, and that rule is the inner function. (2.7.C.1)
\end{itemize} \\
\textbf{The distinction the lesson is built on.} &
\begin{itemize}[leftmargin=1.2em,itemsep=2pt,topsep=0pt]
  \item Composition is \textbf{not} commutative: swapping the machines swaps the
        handoff, so $f(g(x))$ and $g(f(x))$ are typically different functions.
        (2.7.A.3)
  \item \textbf{Target misconception:} students treat $f(g(x))$ and $g(f(x))$ as
        the same thing and compose in whichever order is easier to write ---
        ``same two rules, so same answer.'' Problem 12 of the notes runs one
        input through both orders off one pair of functions and gets 8 and 14.
  \item A near neighbour worth naming out loud: $(f \circ g)(x)$ is not
        $f(x) \cdot g(x)$. That was last lesson's operation, and a product has no
        handoff at all.
\end{itemize} \\
\end{tabularx}
\end{skillbox}

\begin{skillbox}[{Vocabulary, Concepts, \& Theorems}]{greenbox}
\renewcommand{\arraystretch}{1.4}
\begin{tabularx}{\linewidth}{@{} >{\leavevmode\bfseries\color{plum}}p{0.27\linewidth} X @{}}
Composition            & Running two functions back to back. Written $(f \circ g)(x) = f(g(x))$ and read ``$f$ of $g$ of $x$.'' The circle is not multiplication. \\
Inner and outer function & In $f(g(x))$, the \textit{inner} function $g$ runs first; the \textit{outer} function $f$ runs on its result. The letter written last runs first. \\
Handoff value          & The number (or expression) the inner function passes over: its \textit{answer}, and immediately the outer function's \textit{question}. \\
Domain of a composite  & The inputs $x$ that $g$ accepts \textit{and} whose handoff $g(x)$ is a value $f$ accepts. A handoff the outer function refuses gives no value. \\
Decomposition          & Writing a single function $h$ as $f(g(x))$ for simpler pieces --- composition run backwards, done by naming the handoff. \\
\end{tabularx}
\end{skillbox}

\begin{skillbox}[Activate Prior Knowledge \& Spiral Review (5 min)]{lilac}
    Please see attached warm-up (spiral review).
    Skills reviewed: evaluating $f(x) = 3x - 2$ at a value and at a negative value
    (1.1), and at an \textit{expression}, $f(x^2)$ --- which is exactly the
    substitution Part 3 needs, done once before it has a name;
    evaluating a piecewise function by deciding which rule owns the input (1.1);
    adding and subtracting two functions from a table, $(f+g)(2)$ and $(f-g)(4)$
    (1.2) --- the operation students will confuse with composition all period;
    and reading a table \textit{backwards}, which is the habit Part 2 needs when a
    handoff must be found as an input.
\end{skillbox}

\begin{skillbox}[Hook (2 min)]{lilac}
    \textbf{Two registers, one jacket.} A \$60 jacket, a \$10 coupon, 6\% sales
    tax. Register 1 scans the coupon first, then the tax; register 2 scans the
    tax first, then the coupon.
    \begin{enumerate}[itemsep=2pt]
        \item Take a vote before any computing: same total, or different? Most
              classes split, and most students argue ``same --- it's the same two
              things.''
        \item Then work both on the board, and \textbf{write the middle number on
              the board too}: register 1's screen reads \$50 before it charges
              \$53.00; register 2's screen reads \$63.60 before it charges
              \$53.60.
    \end{enumerate}
    \textit{Discussion prompt:} the totals differ because the \textbf{middle
    numbers} differ. That number on the little screen is today's whole idea: it
    is what the first function answered, and it is what the second function gets
    asked. Name it the \textbf{handoff} and keep the word all period.
\end{skillbox}

\begin{skillbox}[Lesson --- I Do / We Do / You Do]{lilac}
\begin{multicols}{2}

\textbf{I DO --- Part 1: The Handoff, and a Table} (10 min)

\textit{Notes rows THE HANDOFF and FINDING IT: read each definition, mark up the
display, work problems 1 and 3. Teacher works; students watch and annotate.
Nothing is cold-called in this phase.}
\begin{itemize}
  \item Read the notation once and insist on it: $(f \circ g)(x) = f(g(x))$, read
        ``$f$ of $g$ of $x$.'' Say out loud what it is \emph{not} --- the circle
        is not multiplication, and $f(x)\cdot g(x)$ was last lesson.
  \item Mark up the two-machine display yourself. Trace the gold arrow with a
        finger and label it twice: ``$E$'s answer'' above, ``$T$'s question''
        below. Fill the two label boxes --- \textit{feet}, \textit{question}.
  \item The spine, repeated verbatim all period: \textbf{one number, two jobs ---
        an answer first, then a question. Find it, and write it down.}
  \item Work \textbf{problem 1} aloud on the cable car: the handoff is
        $E(4) = 800$ \emph{feet}, and only then $T(800) = 60$ degrees. Say the
        units; they are why the order cannot be reversed.
  \item Read the FINDING IT sentences, then work \textbf{problem 3} in two
        recorded steps --- handoff $g(1) = 3$, then $f(3) = 2$. Never skip the
        first step, even though you can do it in your head.
\end{itemize}

\textbf{WE DO --- Part 2: Both Orders in a Table} (6 min)

\textit{Problems 2 and 4--6, and the handoff row of the small table. Class fills
the blanks together; cold-call every one.}
\begin{itemize}
  \item Make the class fill the \textbf{handoff row} all the way across before
        anyone looks up a single value of $f$. That row is the habit.
  \item \textbf{Problem 4} is problem 3 with the machines swapped: handoff 4,
        answer 4. Ask what changed --- the handoff, before anything else did.
  \item \textbf{Problem 5} is the domain of a composite at exposure depth: the
        handoff is 7 and the table has no $f(7)$, so there is no value. No
        interval algebra today.
  \item \textbf{Problem 2} lands the language: $E$ runs first, and
        $(T \circ E)(t)$ does not mean $T(t)\cdot E(t)$.
\end{itemize}

\columnbreak

\textbf{WE DO --- Part 3: Building the Composite} (7 min)

\textit{Notes row BUILDING IT. Class fills the blanks together; cold-call every
one. Teacher works problem 7; the class works 8--10.}
\begin{itemize}
  \item One sentence of method, then the steps: write $f$ with an empty slot, put
        $g(x)$ in \emph{every} slot, simplify. Say ``every,'' and have students
        circle each $x$ before replacing anything.
  \item Point back at the warm-up: they already substituted $x^2$ into
        $3x - 2$. The only new thing is that the substituted expression has a
        name --- it is the handoff.
  \item \textbf{Problem 8} produces $9x^2 - 12x + 5$ against problem 7's
        $3x^2 + 1$: not even the same \emph{shape}. That is the algebraic form of
        the hook.
  \item \textbf{Problem 9} is the ``two routes, one number'' check: handoff 5,
        then $f(5) = 13$; or straight from $3x^2+1$. Same 13 either way.
\end{itemize}

\textbf{WE DO --- Part 4: Two Handoffs, and Taking One Apart} (8 min)

\textit{Notes row TWO HANDOFFS --- the spine of the lesson. Class fills the
blanks together; cold-call every one. Teacher works problem 11; the class works
12--14.}
\begin{itemize}
  \item Mark up the register display before any algebra: two rows, same \$60 in,
        two boxed middle numbers, two different totals. Fill the label boxes ---
        \$0.60, and \textit{handoffs} --- then give one minute for the ``in your
        own words'' line.
  \item Work \textbf{problem 11}: $1.06p - 10.60$ against $1.06p - 10$. The gap
        is \$0.60 at \emph{every} price, because taxing first taxes the ten
        dollars you never paid.
  \item \textbf{Problem 12 is the crux}: one input, 10, through both orders of
        the same pair --- handoff 4 gives 8, handoff 20 gives 14. Ask for the two
        handoffs before either answer.
  \item Read the decomposition paragraph, then \textbf{problem 13}: what would a
        calculator do first in $\sqrt{x+5}$? That rule is the inner function.
        Make them substitute back. \textbf{Problem 14} asks for the
        misconception in prose.
\end{itemize}

\textbf{YOU DO --- Part 5: The You Do Row} (14 min)

\textit{Notes row TABLE \& FORMULAS, problems 15--18 --- the whole independent
phase. Students work alone; circulate and say nothing a neighbor could say
instead.}
\begin{itemize}
  \item \textbf{Problem 15} is the crux transferred to a table: $p(q(2)) = 5$ but
        $q(p(2)) = 2$, with no formula to hide behind.
  \item \textbf{Problem 16} builds both orders for $x+3$ and $x^2$ ---
        $x^2 + 3$ against $x^2 + 6x + 9$; \textbf{17} checks one value of it two
        ways, at a negative input.
  \item \textbf{Problem 18} is the finisher's question: decompose $(2x+1)^3$,
        say which rule runs first, and evaluate at $x = 1$. Naming the order is
        the part that is hard, and homework problem 8 asks it again.
\end{itemize}
\end{multicols}
\end{skillbox}

\begin{skillbox}[Active Monitoring]{redbox}
Circulate during the \textbf{You Do} phase --- the fourteen minutes where students are working
problems 15--18 without you --- and check. The crux is \textbf{problem 12} ($f(g(10)) = 8$ but
$g(f(10)) = 14$); its You Do transfer is \textbf{problem 15} (the same clash, read off a table).
\begin{itemize}
  \item \textbf{Common error (the crux):} composing in whichever order is easier to write, on
        the theory that the same two rules must give the same answer. Cold-call: ``Which function
        is inside? What number does it hand over? Now say the other one's handoff.''
  \item \textbf{Common error:} no handoff written down at all. Chase this hardest --- a student
        who does $p(q(2))$ in their head has no place to be wrong out loud and no place for you
        to catch it. Ask for the middle number before the answer, every time.
  \item \textbf{Common error:} reading $(f \circ g)(x)$ as $f(x) \cdot g(x)$ --- last lesson's
        operation. Ask where the handoff is; there isn't one in a product, and that ends it.
  \item \textbf{Common error:} running the outer function first --- computing $p(2)$ when the
        problem says $p(q(2))$. Point at the parentheses: innermost first, like any arithmetic.
  \item \textbf{Common error:} substituting into only the first $x$, so $f(x) = 3x - 2$ becomes
        $3(x^2+1) - 2x$. Have them circle every $x$ before replacing anything.
  \item \textbf{Common error:} decomposing backwards --- calling $f(u) = 2u + 1$ the outer
        function of $(2x+1)^3$. Make them substitute back; they built $2x^3 + 1$ instead.
\end{itemize}
\end{skillbox}

\begin{skillbox}[Differentiation --- During You Do (14 min)]{redbox}
There is no separate activity sheet: students work the notes' You Do row (problems 15--18) alone
while you circulate. Differentiate by \textit{where you stand}, not by handing out different
paper.

\begin{itemize}
  \item \textbf{Support.} Sit with a struggling student on \textbf{problem 15} and make them
        write the two handoffs in a column before either answer --- $q(2) = 1$ over $p(2) = 3$.
        Seen side by side the two middle numbers do the arguing themselves; a student who writes
        only the answers has nothing to compare.
  \item \textbf{On level.} \textbf{Problems 16 and 17} are what the class should clear unaided.
        Do not answer 17: have them say which number the inner function handed over, then check
        it against their own rule from 16. Two routes, one number --- then walk away.
  \item \textbf{Extend.} \textbf{Problem 18} is the finisher's question, and the hard half is
        naming which rule runs \textit{first}, not producing the split. A student who finishes
        early gets one more: ``find a function $g$ so that $f(g(x)) = x$ for $f(x) = 5x - 2$'' ---
        anyone who produces $\dfrac{x+2}{5}$ has invented an inverse a day early and should
        present it in the Debrief.
\end{itemize}
\end{skillbox}

\boxguard[20]
\begin{skillbox}[Debrief (5 min)]{redbox}
Whole class, teacher-led, packets still open. Three moves, in order:
\begin{enumerate}[itemsep=3pt]
  \item \textbf{Share out the You Do} (3 min) --- one answer per problem, not a full review.
        Problem 15: both handoffs, then both answers. Problem 16: the two rules side by side.
        Problem 17: the value, and which route the student took. Problem 18: the split, and which
        rule runs first. If a student found a $g$ with $f(g(x)) = x$ during the extension, they
        present it here --- it is the on-ramp to 1.4 and the only preview the rest of the class
        gets.
  \item \textbf{Name the headline} (1 min) --- say it, then make the class say it back:
        \textit{one number sits between the two machines --- the inner function's answer and the
        outer function's question. Swap the machines and you swap the handoff.} Point at the two
        boxed middle numbers on the register display while you say it.
  \item \textbf{Point forward} (1 min) --- 1.4 hunts for the function that \emph{undoes} another,
        and the test for it is a composition run in both directions. Mention, without teaching
        it, that in 1.5 a horizontal shift will turn out to be a composition too.
\end{enumerate}
\textbf{Do not} re-teach the substitution method here, take new questions, or let the homework
start early. If the launch shows it did not land, it is the 1.4 warm-up's job, not this minute's.
\end{skillbox}

\begin{skillbox}[Homework Launch (3 min)]{redbox}
\textbf{There is no exit ticket.} The formative read comes twice: circulating the You Do, and
again here. Write the due date on the board, have students copy it into the cover's Homework row,
then start \textbf{homework problems 1 and 2} alone and in silence (3 minutes). Circulate straight
to the students you flagged during You Do.
\begin{itemize}[itemsep=2pt]
  \item \textbf{Problem 1(a) and 1(b)} are the habit: a handoff blank filled before an answer
        blank. A student who fills only the answers has not formed it.
  \item \textbf{Problem 2 is the diagnostic}: the same pair of rules, the single input 8, both
        orders --- handoff 3 gives 12, handoff 32 gives 27. Two matching answers in 2(a) and 2(b)
        means the student ran the same machine twice; a correct pair of answers with both handoff
        blanks empty is a different diagnosis, and only the first needs the crux retaught.
\end{itemize}
\textbf{Sort what you see into three piles} as you walk: \textit{order solid} (nothing to do),
\textit{order shaky} (a 1.4 warm-up item running one input through both orders), and \textit{still
multiplying $f$ by $g$} (a one-minute check-in at the door of 1.4).
\end{skillbox}

\begin{skillbox}[Reinforcement \& Extension]{goldbox}
\textbf{Homework --- printed, graded (2 pages).} Last in the packet; due on the date the teacher
sets. Every context is new. \textit{Part A} (1--3) works handoffs off a table, including one
composite with no value (1c) and the \textbf{diagnostic}, problem 2. \textit{Part B} (4--5) builds
both composites for two pairs of formulas and checks one value two ways. \textit{Part C} (6--7) is
composition in context: a step counter chained to a calorie rule, and an ounces-to-pounds
conversion chained to a price, where the reversed order has the wrong \emph{units}. \textit{Part
D} (8--10) decomposes two functions and checks by substituting back, meets a pair that undoes
itself (9 --- the on-ramp to 1.4), and problem 10 asks for the headline in writing.

\textbf{AP Practice --- extra credit (2 pages).} Optional, before the homework in the packet.
Section I is five AP-style multiple-choice items (four options) whose distractors are the lesson's
errors: composing in the wrong order (1A, 2C, 3A, 5C), multiplying $f(x)g(x)$ instead of composing
(1B, 2D, 3C, 5D), stopping at the handoff and reporting it as the answer (1D, 2A, 5A), and
decomposing backwards (4B). Section II is one free-response question on a wind farm given by two
tables: the handoff with units and an interpretation (a), why the reversed order is meaningless
(b), two inputs sharing one handoff (c), a handoff the outer function will not accept (d), and a
three-function chain (e) --- the finisher's question.

\textbf{Preview:} Next lesson (1.4, Inverse Functions) answers the question the You Do extension
raises: which function \emph{undoes} $f$? The test is a composition run in both directions,
$f(f^{-1}(x)) = x$ and $f^{-1}(f(x)) = x$, so today's habit of naming the handoff is exactly the
tool. Later, 1.5 will reveal a horizontal shift to be a composition as well --- do not mention
that as something already taught.
\end{skillbox}

\begin{teachernote}[Warm-Up]
Five minutes, silent start, no calculators. Item 1(b) is the payoff: substituting $x^2$ into
$f(x) = 3x - 2$ is a routine algebra question here and the entire mechanic of Part 3 --- when you
get there, say out loud that they already did it, and that the thing they substituted was the
handoff. Item 1(a) is not filler either: $f(5) = 13$ is the exact number Part 3 lands on when it
checks a value two ways, and students who recognize it get the ``two routes, one number'' point
for free. Item 2 keeps 1.1's piecewise decision warm. Item 3(a) is the one to watch: $(f+g)(2)$ is
last lesson's operation, and it is the thing students will confuse with composition within twenty
minutes --- do it here so the contrast is available later. Item 3(b) reads the table backwards,
which is the habit Part 2 needs when a handoff has to be found as an \textit{input}.
\end{teachernote}

\begin{teachernote}[Guided Notes]
The notes are one Main Ideas / Notes table, four instruction rows and a You Do row; in each
instruction row you read the printed definition, mark up the display, and work the first problem
of the grid (1, 3, 7, 11), and the class works the rest with the pen in their hand, every problem
cold-called.

The single most valuable thing you can enforce today is \textbf{writing the handoff down}, and the
notes are built to make that unavoidable: the two-step wording in the FINDING IT row, the
handoff row of the small table that must be filled all the way across before a single value of $f$
is looked up, and the two boxed middle numbers in the register display. Do not let a student skip
the first step because they can do it in their head --- the habit, not the arithmetic, is what
carries into 1.4 and into Unit 4.

TWO HANDOFFS is the spine. Do the display before the algebra: the same \$60 entering two rows and
leaving as two different totals, with the two handoffs boxed on the dashed line, is the
misconception drawn. Then problem 12 asks for the two handoffs before either answer: 4 and 20,
which is \emph{why} 8 and 14. Do not accept ``order matters'' as a slogan; make them name the
middle numbers. Problem 13 turns the same idea around into decomposition, and problem 14 asks for
it in prose.

Domain of a composite stays at exposure depth --- the table value $f(g(2))$ with no answer is
enough, and no student needs interval algebra today. The You Do row (15--18) is the entire
independent phase and carries fourteen minutes, so treat it as the lesson's third act. If the
clock slips, cut problem 10 and let the Debrief carry it.
\end{teachernote}

\begin{teachernote}[Debrief]
Five minutes, and on this lesson they are genuinely at risk --- the guided phase runs 21 minutes
and will try to borrow them. Borrow from problem 10 instead; it is the one piece of the period
that is a second example rather than a new idea. Do not borrow from the three-minute Homework
Launch either: with no exit ticket, it is the last look you get before the work leaves the room.
The share-out is one answer per You Do problem so that a student who stalled on 15 still hears the
decomposition result in 18 --- the exact shape of homework problem 8, which they will otherwise
meet alone at home. Say the headline yourself first and make the class say it back; the sentence
they need is about the middle number, not about order.
\end{teachernote}

\begin{teachernote}[AP Practice]
Extra credit, optional, and not a substitute for the homework --- say so when you hand out the
packet. Score it however extra credit works in your gradebook; the cover gives it a $+$ cell
outside the total. Multiple-choice answers are \textbf{C, B, D, A, B}. Item 2 is the most useful
one to glance at: (C) is the wrong order, (A) is the handoff reported as the answer, and (D) is
$f(3) \cdot g(3)$ --- three different diagnoses on one line, and any of them means the crux did not
land. In the free response, part (b) is the units argument and part (d) is the domain of a
composite; both should be written in context. Part (c) is the subtle one --- two different times
produce the same handoff, so they must produce the same power; credit any answer that says so.
\end{teachernote}

\begin{teachernote}[Homework]
Graded, two pages, due on the date you set; students start problems 1 and 2 in the Homework Launch.
Grade problems 2 and 10 for accuracy and the rest for completion and shown work if time is short:
2 (both orders off one input) and 10 (the headline in writing) are the two items that say whether
the crux landed, and together they predict 1.4 better than the building items do. Problem 1(c) is
the only place the domain of a composite is written in prose --- scan it alongside 7(b), which asks
the same question with units instead of a table. Problem 9 is the on-ramp to 1.4: a student who
notices that both composites come out to $x$ has met the inverse relationship before it has a name,
and their answer is the best possible opening for tomorrow. On problem 8, a student who writes
inner $u^5$ and outer $4x-1$ has decomposed backwards; have them substitute back and watch it
rebuild the wrong function.
\end{teachernote}
\end{document}
